% Part: sets-functions-relations
% Chapter: arithmetization
% Section: cauchy

\documentclass[../../../include/open-logic-section]{subfiles}

\begin{document}

\olfileid{sfr}{arith}{cauchy}

\olsection{అనుబంధం: కౌషీ క్రమాలుగా వాస్తవ సంఖ్యలు}

\olref[cuts]{sec}లో వాస్తవ సంఖ్యలను డెడెకిండ్ కోతలుగా నిర్మించాం.
ఈ విభాగంలో ఒక ప్రత్యామ్నాయ నిర్మాణాన్ని వివరిస్తాం. ఇది కౌషీ ఇచ్చిన
(ఇప్పుడు మనం పిలిచే) కౌషీ క్రమ నిర్వచనంపై ఆధారపడుతుంది; అయితే ఈ
నిర్వచనాన్ని వాస్తవ సంఖ్యలను \emph{నిర్మించడానికి} ఉపయోగించిన ఘనత
పంతొమ్మిదవ శతాబ్దపు ఇతర రచయితలకు, ముఖ్యంగా వైయర్‌స్ట్రాస్, హైనె,
మేరే, కాంటర్‌లకు చెందుతుంది. (మంచి చరిత్ర కోసం
\citeauthor{OConnorRobertson:RN} \citeyear{OConnorRobertson:RN} చూడండి.)

పంతొమ్మిదవ శతాబ్దానికి వెళ్లేముందు సైమన్ స్టెవిన్ (1548--1620)ను
పరిశీలించడం ఉపయోగకరం. సంక్షిప్తంగా, ప్రతి వాస్తవ సంఖ్యను దాని దశాంశ
విస్తరణ రూపంలో భావించవచ్చని స్టెవిన్ గ్రహించాడు. అందువల్ల $\sqrt{2}$
లాంటి అకరణీయ సంఖ్యకు కూడా చక్కని దశాంశ విస్తరణ ఉంది; దాని ఆరంభం:
\[
	1.41421356237\ldots
\]
సమితి సిద్ధాంతంలో దశాంశ విస్తరణలను నమూనీకరించడం చాలా సులభం: వాటిని
$d \colon \Nat \to \Nat$ అనే ప్రమేయాలుగా పరిగణిస్తే చాలు; ఇక్కడ
$d(n)$ మనకు కావలసిన $n$వ దశాంశ స్థానం. దశాంశ బిందువుకు ముందు ఉండే
వాస్తవ సంఖ్య భాగాన్ని (ఇక్కడ కేవలం $1$) నిర్వహించడానికి కొద్దిగా
సవరణ కావాలి. ఉదాహరణకు $0.999\ldots = 1$ అని హామీ ఇవ్వడానికి మరో
సవరణ (ఒక తుల్యతా సంబంధం) కూడా కావాలి. అయినా స్టెవిన్ పద్ధతిలో సమితి
సిద్ధాంతం లోపల వాస్తవ సంఖ్యలకు పూర్తిగా కచ్చితమైన నిర్మాణాన్ని ఇవ్వడం
కష్టం కాదు.

మన దృష్టి స్టెవిన్‌పై కాదు. (స్టెవిన్ గురించి మరింత సమాచారం కోసం
\citealt{KatzKatz2012} చూడండి.) కానీ దీనికి దగ్గరగా ఉన్న ఆలోచన ఇదిగో.
$\sqrt{2}$ దశాంశ విస్తరణను నేరుగా తీసుకునే బదులు, మరింత కచ్చితమైన
విస్తరణలను పరిశీలిస్తూ $\sqrt{2}$కు క్రమంగా కచ్చితమయ్యే కరణీయ
ఉజ్జాయింపుల \emph{క్రమాన్ని} తీసుకోవచ్చు:
\[
1, 1.4, 1.414, 1.4142, 1.41421,\ldots
\]
వాస్తవ సంఖ్యలను ``క్రమంగా మెరుగయ్యే ఉజ్జాయింపుల'' ద్వారా
పరిగణించవచ్చనే ఆలోచన, మరో అంతర్దృష్టుల వరుసకు ఆధారం ఇస్తుంది
($\Int$ నుంచి $\Rat$ను లేదా $\Nat$ నుంచి $\Int$ను నిర్మించినప్పుడు
ఉపయోగించిన అవగాహనలను ఇది పోలి ఉంటుంది). ప్రాథమిక అవగాహనలు ఇవి:
\begin{enumerate}
	\item ప్రతి వాస్తవ సంఖ్యను ఒక (బహుశా అనంతమైన) దశాంశ విస్తరణగా
	రాయవచ్చు.
	\item ఒక (బహుశా అనంతమైన) దశాంశ విస్తరణలో సంకేతీకరించిన సమాచారాన్ని
	కరణీయ సంఖ్యల క్రమం ద్వారా కూడా అంతే విధంగా సంకేతీకరించవచ్చు.
	\item కరణీయ సంఖ్యల క్రమాన్ని $\Nat$ నుంచి $\Rat$కు ఒక ప్రమేయంగా
	భావించవచ్చు; క్రమంలోని $n$వ కరణీయ సంఖ్యను $f(n)$గా తీసుకుంటే చాలు.
\end{enumerate}
అయితే $\Nat$ నుంచి $\Rat$కు వెళ్లే \emph{ఏ} ప్రమేయమైనా మనకు వాస్తవ
సంఖ్యను ఇవ్వదు. ఉదాహరణకు ఈ ప్రమేయాన్ని పరిశీలించండి:
\[
	f(n) =\begin{cases}
			1 & \text{if }n\text{ is odd}\\
			0 &\text{if }n\text{ is even}
		\end{cases}
\]
ఇక్కడి ప్రధాన సమస్య ఏమిటంటే, $0,1,0,1,0,1,0,\ldots$ అనే క్రమం ఏ
వాస్తవ సంఖ్యవైపూ ``సమీపిస్తూ'' ఉన్నట్లు కనిపించదు. అందువల్ల ఏదో
వాస్తవ సంఖ్యవైపు సమీపించే క్రమాలనే పరిశీలించేందుకు, ఏదో
\emph{పరిమితి} గల క్రమాలకే మన దృష్టిని పరిమితం చేయాలి.

పరిమితి అనే ఆలోచనను \olref[his][set][limits]{sec}లో ఇప్పటికే
చూశాం. కానీ అక్కడ వాడిన నిర్వచనాన్నే \emph{యథాతథంగా} ఇక్కడ
వాడలేం. అక్కడి ``$(\forall \epsilon>0)$'' అనే వ్యక్తీకరణలో వాస్తవ
సంఖ్యలపై పరిమాణీకరణ అంతర్నిహితంగా ఉంది; అంతేకాక వాస్తవ సంఖ్యలపై
ప్రమేయాల పరిమితులను పరిశీలించాం. అందువల్ల ఆ నిర్వచనాన్ని వాడితే
వాస్తవ సంఖ్యలను ముందే తీసుకున్నట్లవుతుంది; కానీ మనం
\emph{నిర్మించాలని} ప్రయత్నిస్తున్నవి అవే. అదృష్టవశాత్తూ, పరిమితికి
దగ్గరగా ఉన్న ఒక భావనతో పనిచేయవచ్చు.
\begin{defn}\ollabel{def:CauchySequence}
  $f: \Nat \to \Rat$ అనే ప్రమేయం, ప్రతి ధన $\epsilon \in \Rat$కు
  $(\exists \ell \in \Nat)(\forall m, n > \ell)|f(m) - f(n)| <
  \epsilon$ అయినప్పుడు మాత్రమే ఒక \emph{కౌషీ క్రమం}.
\end{defn}

పరిమితి గురించిన సాధారణ ఆలోచన మునుపటిదే: మీకు ఒక నిర్దిష్ట కచ్చితత్వ
స్థాయి ($\epsilon$తో కొలిచేది) కావాలంటే, చూడాల్సిన ఒక ``ప్రాంతం''
ఉంటుంది ($\ell$ కంటే పెద్ద ఏ ఇన్‌పుట్ అయినా). మన $1$, $1.4$,
$1.414$, $1.4142$, $1.41421$\ldots క్రమానికి పరిమితి ఉందని సులభంగా
చూడవచ్చు: $\nicefrac{1}{10^{n}}$ లోపం లోపల $\sqrt{2}$ను
ఉజ్జాయించాలంటే, $n$వ స్థానం తరువాతి ఏ పదాన్నైనా చూడండి.

అప్పుడు సహజంగా వచ్చే ఆలోచన, ఏ కౌషీ క్రమమైనా ఒక వాస్తవ సంఖ్యే అని
చెప్పడం. కానీ $\Int$, $\Rat$ నిర్మాణాల్లో మాదిరిగా ఇది మరీ సరళం:
ఇచ్చిన ఒక్క వాస్తవ సంఖ్యను అనేక భిన్నమైన కౌషీ క్రమాలు సూచిస్తాయి.
దీనిని చూడటానికి ఒక సులభమైన మార్గం ఇది. ఒక కౌషీ క్రమం $f$ ఇచ్చినప్పుడు,
$g(0)\neq f(0)$ కావడం తప్ప $f$తో పూర్తిగా సమానంగా ఉండే ప్రమేయంగా
$g$ను నిర్వచించండి. మొదటి సంఖ్య తరువాత అన్ని చోట్లా రెండు క్రమాలు
ఏకీభవిస్తాయి కనుక, \olref{def:CauchySequence}లో ఉపయోగించిన అర్థంలో
వాటికి ఒకే పరిమితి ఉందని (చివరికి) చెప్పాలనుకుంటాం; అందువల్ల అవి ఒకే
వాస్తవ సంఖ్యను ``నిర్వచిస్తున్నట్లు'' భావించాలి. కాబట్టి ఈ కౌషీ
క్రమాలను నిజంగా ఒకే వాస్తవ సంఖ్యగా భావించాలి.

అందువల్ల కౌషీ క్రమాలపై మళ్లీ ఒక తుల్యతా సంబంధాన్ని నిర్వచించి,
వాస్తవ సంఖ్యలను తుల్యతా వర్గాలతో గుర్తించాలి. మొదట పరిమితిలో $0$కు
అభిసరించే ప్రమేయం అనే భావన కావాలి. ఏ $h : \Nat \to \Rat$కైనా,
ప్రతి ధన $\epsilon \in \Rat$కు
$(\exists \ell \in \Nat)(\forall n > \ell)|h(n)| < \epsilon$
అయితే మాత్రమే \emph{$h$ $0$కు అభిసరిస్తుంది} అని అందాం.\footnote{దీనిని
\olref[his][set][limits]{sec}లోని
$\lim_{x \mathord{\rightarrow}\infty}f(x) = 0$ నిర్వచనంతో పోల్చండి.}
ఇంకా $f$, $g$లు $\Nat \to \Rat$ ప్రమేయాలు అయినప్పుడు,
$(f-g)(n) = f(n) - g(n)$గా తీసుకుందాం. ఇప్పుడు ఇలా నిర్వచించండి:
\[
	f \Realequiv g \text{ iff $(f-g)$ tends to $0$}.
\]
$\Realequiv$ ఒక తుల్యతా సంబంధమని తనిఖీ చేయాలి; అది నిజమే. అప్పుడు
మనకు ఇష్టమైతే, $\Nat \to \Rat$ కౌషీ క్రమాలన్నిటి $\Realequiv$
కింద తుల్యతా వర్గాలే వాస్తవ సంఖ్యలు అని నిర్వచించవచ్చు.
\sourcecorrection{OLTEARITH-005}{ఆంగ్ల మూలం ఇక్కడ వాస్తవ సంఖ్యలను
``తుల్యతా సంబంధాలతో'' గుర్తించాలని చెప్పింది; కానీ వెంటనే వచ్చే
నిర్వచనం, మొత్తం నిర్మాణం వాటిని ఆ సంబంధం కింద ఉన్న
\emph{తుల్యతా వర్గాలతో} గుర్తిస్తాయి. ఇక్కడ వస్తువు రకాన్ని తుల్యతా
వర్గాలుగా సరిచేశాం.}

\begin{prob}
ప్రతి $n$కు $f(n) = 0$గా తీసుకోండి. $g(n) = \frac{1}{(n+1)^2}$గా
తీసుకోండి. రెండూ కౌషీ క్రమాలని, నిజానికి రెండు ప్రమేయాల పరిమితీ
$0$ అని, అందువల్ల $f \sim_\Real g$ కూడా అని చూపండి.
\end{prob}

ఇది చేసిన తరువాత, మామూలుగానే $f$ను \tetoken{ఒక మూలకంగా}{!!a{element}} కలిగిన తుల్యతా
వర్గాన్ని $\equivrep{f}{\Realequiv}$గా రాస్తాం. అయితే చదవడానికి
సులభంగా, ఇక నుంచి అధోలిపిని వదిలి కేవలం
$\equivrep{f}{}$ అని రాస్తాం. ఇంకా ప్రతి $q \in \Rat$కు
$q_{\Real} = \equivrep{c_{q}}{}$ అని నిర్దేశిస్తాం; ఇక్కడ ప్రతి
$n \in \Nat$కు $c_q(n) = q$ అయ్యే స్థిర ప్రమేయమే $c_{q}$. తరువాత
వాస్తవ సంఖ్యలపై ప్రాథమిక సంబంధాలు, క్రియలను, ఉదాహరణకు, ఇలా
నిర్వచిస్తాం:
\begin{align*}
	\equivrep{f}{} + \equivrep{g}{} &= 	\equivrep{(f + g)}{} \\
	\equivrep{f}{} \times 	\equivrep{g}{} &= \equivrep{(f \times g)}{}
\end{align*}
ఇక్కడ $(f + g)(n) = f(n) + g(n)$ మరియు
$(f \times g)(n) = f(n) \times g(n)$. $f$, $g$లు కౌషీ క్రమాలు
అయినప్పుడు $(f + g)$, $(f-g)$, $(f\times g)$లలో ప్రతి ఒక్కటీ కౌషీ
క్రమమేనని కూడా తనిఖీ చేయాలి; అవి అలాగే ఉంటాయి, నిరూపణను మీకే
వదిలేస్తాం.

చివరగా ఒక క్రమ భావనను నిర్వచిస్తాం. $\equivrep{f}{}\neq 0_\Real$
మరియు $(\exists \ell \in \Nat)(\forall n > \ell)0 < f(n)$ రెండూ
వర్తిస్తే మాత్రమే $\equivrep{f}{}$ \emph{ధనాత్మకం} అని అందాం.
తరువాత $\equivrep{(g - f)}{}$ ధనాత్మకం అయితే మాత్రమే
$\equivrep{f}{} < \equivrep{g}{}$ అని అందాం. ఇది సునిర్వచితమని
(అంటే తుల్యతా వర్గం నుంచి ఎంచుకున్న ``ప్రాతినిధ్య'' ప్రమేయంపై ఇది
ఆధారపడదని) తనిఖీ చేయాలి.
\sourcecorrection{OLTEARITH-006}{ఆంగ్ల మూలం వాస్తవ సంఖ్యను సూచించే
$\equivrep{f}{}$ను $0_\Rat$తో పోల్చింది. ముందే నిర్వచించిన కరణీయ
సున్నా యొక్క వాస్తవ అంతఃస్థాపనకు సరిపోయేలా దానిని $0_\Real$గా
సరిచేశాం.}
%\begin{align*}
%	\equivrep{f}{} \leq \equivrep{g}{} \text{ iff }&\equivrep{f}{} = \equivrep{g}{} \text{ or }\\
%	&(\exists \ell \in \Nat)(\forall n \in \Nat)(\ell < n \lif f(n) \leq g(n))
%\end{align*}
ఇది చేసిన తరువాత, ఇవి సరైన బీజగణిత ధర్మాలను ఇస్తాయని చూపడం చాలా
సులభం; అంటే:
\begin{thm}\ollabel{thm:cauchyorderedfield}
కౌషీ క్రమాల తుల్యతా వర్గాలు ఒక క్రమిత క్షేత్రాన్ని ఏర్పరుస్తాయి.
\end{thm}
\sourcecorrection{OLTEARITH-007}{ఆంగ్ల మూలంలోని సిద్ధాంతం, సాధన ముడి
కౌషీ క్రమాలే క్రమిత క్షేత్రాన్ని ఏర్పరుస్తాయని చెబుతాయి. కానీ సమానత్వం,
క్రియలు $\Realequiv$ కింద తుల్యతా వర్గాలపై నిర్వచించబడ్డాయి. రెండు
ప్రకటనలలోనూ క్షేత్రపు వస్తువులను ఆ తుల్యతా వర్గాలుగా సరిచేశాం.}

\begin{proof}
సాధన.
\end{proof}

\begin{prob}
కౌషీ క్రమాల తుల్యతా వర్గాలు ఒక క్రమిత క్షేత్రాన్ని ఏర్పరుస్తాయని
నిరూపించండి.
\end{prob}

ఇలా నిర్మించిన వాస్తవ సంఖ్యలకు సంపూర్ణతా ధర్మం ఉందని నిరూపించడం
మరింత కష్టం; అందువల్ల ఆ నిరూపణను ఇస్తాం.

\begin{thm}
శూన్యేతర కౌషీ క్రమాల సమితి ప్రాతినిధ్యం చేసే తుల్యతా వర్గాలకు ఒక
ఊర్ధ్వ అవధి ఉంటే, ఆ వర్గాలకు కనిష్ఠ ఊర్ధ్వ అవధి ఉంటుంది.
\end{thm}
\sourcecorrection{OLTEARITH-008}{ఆంగ్ల మూలం ఈ ప్రకటనలో $S$ను కౌషీ
క్రమాల సమితిగా పిలిచి, నిరూపణలో కొన్నిచోట్ల దాని మూలకాలను నేరుగా
వాస్తవ సంఖ్యలుగా, మరికొన్నిచోట్ల వాటి తుల్యతా వర్గాలుగా పోల్చింది.
ఇక్కడ ఫలితాన్ని తుల్యతా వర్గాల కుటుంబానికి ప్రకటించి, నిరూపణలో $S$ను
ఆ వర్గాల ప్రాతినిధ్య క్రమాల కుటుంబంగా నిలిపి, ప్రతి పోలికను వర్గంపై
స్పష్టంగా రాశాం.}

\begin{proof}[నిరూపణ రూపురేఖ] తుల్యతా వర్గాలకు ఊర్ధ్వ అవధి ఉన్న
శూన్యేతర ప్రాతినిధ్య కౌషీ క్రమాల కుటుంబం $S$ అనుకుందాం. అప్పుడు
$p_{\Real}$, $S$ ప్రాతినిధ్యం చేసే వర్గాలకు ఊర్ధ్వ అవధి అయ్యే ఏదో
$p \in \Rat$ ఉంది.
$r \in S$ తీసుకోండి; అప్పుడు
$q_{\Real} < \equivrep{r}{}$ అయ్యే ఏదో $q \in \Rat$ ఉంది. కాబట్టి
$S$ ప్రాతినిధ్యం చేసే వర్గాలకు కనిష్ఠ ఊర్ధ్వ అవధి ఉంటే, అది
$q_\Real$, $p_\Real$ల మధ్య (రెండింటితో సహా) ఉంటుంది.

కనిష్ఠ ఊర్ధ్వ అవధిని దిగువ, ఎగువ రెండువైపుల నుంచి ఒకేసారి సమీపిస్తూ
దానిని క్రమంగా గుర్తిస్తాం. ప్రత్యేకంగా, $f$ ఎగువ నుంచి, $g$ దిగువ
నుంచి కనిష్ఠ ఊర్ధ్వ అవధిని సమీపించే లక్ష్యంతో రెండు ప్రమేయాలు
$f, g \colon \Nat \to \Rat$లను నిర్వచిస్తాం. మొదట ఇలా నిర్వచిస్తాం:
\begin{align*}
	f(0) &= p \\
	g(0) &= q
\end{align*}
తరువాత $a_n = \frac{f(n) + g(n)}{2}$ అయినప్పుడు ఇలా తీసుకోండి:\footnote{ఇది
పునరావృత్త నిర్వచనం. కానీ పునరావృత్త నిర్వచనాలు సమంజసమైనవని మనం
\emph{ఇంకా} ఏ కారణమూ ఇవ్వలేదు.}
\begin{align*}
	f(n+1) &=
	\begin{cases}
		a_n &\text{if }(\forall h \in S)\equivrep{h}{} \leq (a_n)_\Real\\
		f(n)&\text{otherwise}
	\end{cases}\\
	g(n+1) &=
	\begin{cases}
		a_n &\text{if }(\exists h \in S)\equivrep{h}{} \geq (a_n)_\Real\\
		g(n) &\text{otherwise}
	\end{cases}
\end{align*}
$f$, $g$ రెండూ కౌషీ క్రమాలు. (దీనిని సులభంగానే తనిఖీ చేయవచ్చు;
కానీ సాధనగా వదిలేస్తాం.) $f$, $g$ల మధ్య భేదం ప్రతి మెట్టులో
సగమవుతుంది కనుక $(f-g)$ ప్రమేయం $0$కు అభిసరిస్తుందని గమనించండి.
అందువల్ల $\equivrep{f}{} = \equivrep{g}{}$.

మొదట $\equivrep{f}{}$, $S$ ప్రాతినిధ్యం చేసే వర్గాలకు ఊర్ధ్వ అవధి
అని, అంటే $(\forall h \in S)\equivrep{h}{} \leq \equivrep{f}{}$
అని చూపుతాం. (తర్కంలో \olref{thm:cauchyorderedfield}ను ఉపయోగిస్తాం.)
$h \in S$ తీసుకుని, వైరుధ్య పద్ధతిలో
$\equivrep{f}{} < \equivrep{h}{}$ అనుకుందాం; అప్పుడు
$0_\Real < \equivrep{(h-f)}{}$. $f$ ఏకదిశగా క్షీణించే కౌషీ క్రమం
కనుక, $\equivrep{(c_{f(n)} - f)}{} < \equivrep{(h-f)}{}$ అయ్యే
ఏదో $n \in \Nat$ ఉంది. కాబట్టి:
\[
	(f(n))_\Real = \equivrep{c_{f(n)}}{} < \equivrep{f}{} + \equivrep{(h-f)}{} = \equivrep{h}{},
\]
కానీ నిర్మాణం ప్రకారం $\equivrep{h}{} \leq (f(n))_\Real$; ఇది
వైరుధ్యం.

తరువాత $\equivrep{f}{} = \equivrep{g}{}$, $S$ ప్రాతినిధ్యం చేసే
వర్గాలకు \emph{కనిష్ఠ} ఊర్ధ్వ అవధి అని చూపుతాం. $j$ ఏదైనా కౌషీ క్రమం
అనుకుని $\equivrep{j}{} < \equivrep{g}{}$ అనుకుందాం. పైన మాదిరిగానే
($g$ \emph{వృద్ధిచెందే} క్రమమనే వాస్తవాన్ని వాడుతూ),
$\equivrep{j}{} < (g(n))_\Real$ అయ్యే $n \in \Nat$ ఉంది. కానీ
నిర్మాణం ప్రకారం $(g(n))_\Real \leq \equivrep{h}{}$ అయ్యే
$h \in S$ ఉంది; కాబట్టి $\equivrep{j}{} < \equivrep{h}{}$,
అందువల్ల $\equivrep{j}{}$, $S$ ప్రాతినిధ్యం చేసే వర్గాలకు ఊర్ధ్వ
అవధి కాదు.
\end{proof}

\end{document}
